\documentclass[11pt]{article}
\usepackage[T1]{fontenc}
\usepackage[utf8]{inputenc}
\usepackage{amsmath,amssymb,mathtools}
\usepackage[a4paper,margin=1in]{geometry}
\usepackage{hyperref}

\title{The Remaining Mathematical Obstruction in the Brenier Formalization}
\author{}
\date{}

\begin{document}
\maketitle

\section*{Purpose}
This note explains the exact mathematical issue in the current formalization of Brenier's theorem.
The short version is:
\begin{itemize}
\item the repository already formalizes the chain
  \[
    \text{convex potential} \Longrightarrow \text{subgradient support}
    \Longrightarrow \text{gradient map a.e.} \Longrightarrow \text{graph plan};
  \]
\item the remaining gap is the passage
  \[
    \text{optimal quadratic coupling} \Longrightarrow \text{convex potential},
  \]
  in the full classical generality;
\item the extra hypothesis \texttt{HasRockafellarBound} is not a genuine Brenier assumption, but an
  artifact of the current real-valued Rockafellar construction.
\end{itemize}

\section*{What the code currently proves}
Let \(\pi \in \Pi(\mu,\nu)\) be a coupling on a finite-dimensional real inner product space \(E\).

The repository currently proves the following.

\subsection*{1. If one already has a convex potential, the plan collapses to a map}
Suppose \(\phi : E \to \mathbb R\) is convex and
\[
  \operatorname{supp}(\pi) \subseteq \operatorname{Graph}(\partial \phi).
\]
If \(\mu \ll m\), where \(m\) is an additive Haar measure, then:
\begin{enumerate}
\item \(\phi\) is differentiable \(m\)-almost everywhere;
\item hence \(\phi\) is differentiable \(\mu\)-almost everywhere;
\item at points of differentiability,
  \[
    \partial \phi(x) = \{\nabla \phi(x)\};
  \]
\item therefore \(\pi\) is concentrated on the graph of \(\nabla \phi\);
\item hence
  \[
    \pi = (\mathrm{id}, \nabla \phi)_{\#}\mu,
    \qquad
    (\nabla \phi)_{\#}\mu = \nu.
  \]
\end{enumerate}

This part is formalized.

\subsection*{2. If one already has a real-valued Rockafellar potential, the plan collapses to a map}
The present Rockafellar file constructs a function
\[
  \phi_\Gamma(x) = \sup \mathcal V_\Gamma(x)
\]
for a rooted family of finite-chain values \(\mathcal V_\Gamma(x)\), but only under the hypothesis
that this supremum is finite for every \(x\).

This is encoded by the predicate \texttt{HasRockafellarBound}.

Under that hypothesis, the code proves:
\begin{enumerate}
\item \(\phi_\Gamma : E \to \mathbb R\) is convex;
\item \(\Gamma \subseteq \operatorname{Graph}(\partial \phi_\Gamma)\);
\item therefore, if \(\Gamma = \operatorname{supp}(\pi)\) and \(\mu \ll m\), then
  \[
    \pi = (\mathrm{id}, \nabla \phi_\Gamma)_{\#}\mu.
  \]
\end{enumerate}

\subsection*{3. Final packaged theorem currently available}
The current endpoint is:
\begin{quote}
If \(\pi\) is an optimal quadratic coupling, \(\mu \ll m\), and \(\pi\) satisfies
\texttt{HasRockafellarBound}, then there exists a convex \(\phi\) such that
\[
  \pi = (\mathrm{id}, \nabla \phi)_{\#}\mu,
  \qquad
  (\nabla \phi)_{\#}\mu = \nu,
\]
and the graph plan of \(\nabla \phi\) is optimal.
\end{quote}

\section*{What \texttt{HasRockafellarBound} actually says}
Fix a set \(\Gamma \subseteq E \times E\) and a base point
\[
  (x_0,y_0) \in \Gamma.
\]
The rooted Rockafellar value set at \(x \in E\) is made of all finite-chain expressions
\[
  \sum_{i<n} \langle y_i, x_{i+1} - x_i \rangle + \langle y_n, x - x_n \rangle
\]
over nonempty chains
\[
  [(x_0,y_0), (x_1,y_1), \dots, (x_n,y_n)] \subseteq \Gamma
\]
starting at the chosen base point.

Then \texttt{HasRockafellarBound} says:
\begin{quote}
there exists a base point in \(\operatorname{supp}(\pi)\) such that for every \(x\), the set of all
those chain values is bounded above in \(\mathbb R\).
\end{quote}

Equivalently, the current code assumes that the pointwise supremum defining the Rockafellar
potential is finite everywhere.

\section*{Why this is not part of classical Brenier}
In the standard mathematical proof of Brenier's theorem, one does \emph{not} assume such a bound.

The reason is that one does not start by insisting that the Rockafellar potential be
\(\mathbb R\)-valued everywhere.

Instead, one usually proceeds in one of the following equivalent ways:
\begin{enumerate}
\item define a proper convex potential that is allowed to take the value \(+\infty\);
\item or define the potential on the convex hull of the relevant domain only;
\item or prove finiteness only on the set where it is later needed.
\end{enumerate}

In all of these formulations, the supremum is allowed to exist without a global bounded-above
hypothesis on every chain-value set.

So the mathematical issue is \emph{not} that Brenier truly needs an extra assumption.
The issue is that the present formalization took the convenient route
\[
  \phi_\Gamma : E \to \mathbb R
\]
too early.

\section*{Where the current proof becomes too rigid}
The rigid point is the combination of the following two decisions:
\begin{enumerate}
\item define the Rockafellar potential as an actual real-valued supremum;
\item feed that directly into the real-valued convex and differentiability API.
\end{enumerate}

That gives a clean downstream pipeline, but it forces an upstream finiteness hypothesis.

More precisely:
\begin{itemize}
\item the current Rockafellar theorem proves
  \[
    \Gamma \subseteq \operatorname{Graph}(\partial \phi)
  \]
  only after one knows \(\phi : E \to \mathbb R\);
\item but classical Rockafellar gives first a \emph{proper convex} potential, not necessarily finite
  on all of \(E\);
\item the current \texttt{ConvexAE} file is a theorem about real-valued convex functions
  \(\phi : E \to \mathbb R\), not proper extended-real convex functions.
\end{itemize}

So the obstruction is structural:
\begin{quote}
the present upstream API is too narrow for the classical Rockafellar construction, while the present
downstream API expects exactly that narrow real-valued input.
\end{quote}

\section*{What the missing theorem should conceptually look like}
The ideal mathematical statement is something like:
\begin{quote}
If \(\Gamma \subseteq E \times E\) is cyclically monotone for the pairing, then there exists a
proper convex function \(\phi\) such that
\[
  \Gamma \subseteq \operatorname{Graph}(\partial \phi).
\]
\end{quote}

Then, in the Brenier setting, one proves that \(\phi\) is finite \(\mu\)-almost everywhere on the
source side, and therefore that \(\nabla \phi\) is meaningful \(\mu\)-almost everywhere.

That is the standard textbook route.

\section*{Why the issue is genuinely mathematical, not merely technical}
One might ask whether the current bound can simply be deduced from cyclical monotonicity.

In general, no: cyclical monotonicity alone does not force a globally finite real-valued
Rockafellar supremum on all of \(E\).

What cyclical monotonicity gives is the geometric structure needed for a supporting convex object.
But that object may naturally be proper and extended-valued.

So the distinction is:
\begin{itemize}
\item \emph{classical Rockafellar}: a proper convex potential;
\item \emph{current implementation}: a globally finite real-valued potential.
\end{itemize}

The latter is stronger, and that is exactly where the extra assumption enters.

\section*{What a correct fix would look like}
To remove the gap in the mathematically right way, the code should be refactored along the following
lines.

\subsection*{Option A. Extended-real Rockafellar first}
\begin{enumerate}
\item Replace the current real-valued Rockafellar potential by a proper convex potential allowed to
  take the value \(+\infty\).
\item Prove the Rockafellar containment theorem at that level:
  \[
    \Gamma \subseteq \operatorname{Graph}(\partial \phi).
  \]
\item Then prove a Brenier-specific finiteness theorem showing that in the quadratic optimal
  transport setting, \(\phi\) is finite \(\mu\)-almost everywhere.
\item Finally adapt the downstream gradient-collapse argument to that proper-convex setting.
\end{enumerate}

\subsection*{Option B. Localized real-valued potential}
\begin{enumerate}
\item Restrict the construction to a set on which the supremum is known to be finite.
\item Prove that this suffices for the support of the optimal plan.
\item Run the differentiability and graph-plan argument only on that region.
\end{enumerate}

This is more ad hoc, and less faithful to the standard convex-analytic statement, but may be easier
to formalize.

\section*{Why I stopped before coding this refactor immediately}
The missing step is no longer a short theorem in the current language. It requires choosing a new
mathematical interface.

In particular, one has to decide:
\begin{itemize}
\item what codomain to use for the proper convex potential;
\item how to formulate subgradients for that codomain;
\item how to connect that formulation back to the existing real-valued differentiability theorems.
\end{itemize}

This is why the current code isolates the issue but does not yet resolve it completely.

\section*{The shortest conversation to have next}
The right next design question is:
\begin{quote}
Should the repository move to a proper extended-valued convex potential, or should it localize the
real-valued construction only where the Brenier argument needs finiteness?
\end{quote}

That is the genuine mathematical fork in the road.

\end{document}
